% !TEX root = ../main.tex
% ============================================================================
% Budget: 2.3 pages. The page holds two full-width tables and two or three
% figures at most; anything more goes to Appendix D.
%
% HOW TO WRITE A SUBSECTION HERE
%   1. Start with the finding in one bold sentence, as the midterm results did.
%   2. Give the numbers that support it and point to the table or figure.
%   3. Say what limits the finding (scene count, interval, what was not run).
% Each subsection below has a comment with what it must show and where the
% numbers come from. Replace every \tbd and \todo; delete the comment last.
%
% WHERE THE NUMBERS COME FROM
%   per-image scores of a head   confhead/examples/evaluate_head.py  (CARC: carc-eval.sbatch)
%   flip baseline                confhead/examples/run_flip_baseline.py
%   tables with intervals        confhead/examples/build_results_table.py
%   figures                      confhead/src/evals/plot.py  -> ../figures-src/
%
% BEFORE ANY NUMBER GOES IN (rubric gates 2 and 3)
%   - All methods of a table are scored on the same validation images and cells.
%   - Intervals come from the paired bootstrap over scenes (Section 4.4).
%   - The label masks D8, D9, D13 are settled (Section 3.3).
% ============================================================================
\section{Results}
\label{sec:results}

% ----------------------------------------------------------------------------
\subsection{Main Comparison}
\label{sec:main}
% MUST SHOW: the head against every score that needs no training, and against
% the two bounds, on all validation images. One row per method of Section 4.3.
% The reader should see at a glance (a) how far the best head is from the
% oracle, (b) how much it gains over flip disagreement, and (c) what the
% lidar-free head loses against the full head.
% Negative results stay in the table (rubric gate 5): if a simple score matches
% a head, the row says so.

\todo{One bold sentence with the main finding, then the numbers of Table~\ref{tab:main} and the reading of Figure~\ref{fig:sparsification}.}

\begin{table*}[t]
  \centering
  \small
  \setlength{\tabcolsep}{5pt}
  \begin{tabular}{@{}l l c c c c c c@{}}
    \toprule
    \textbf{Method} & \textbf{Reads} & \textbf{Lidar at inference} & \textbf{AUSE} $\downarrow$ & \textbf{AURG} $\uparrow$ & \textbf{AUROC} $\uparrow$ & \textbf{ECE} $\downarrow$ & \textbf{Brier} $\downarrow$ \\
    \midrule
    Random order & -- & no & \tbd & 0 & 0.5 & -- & -- \\
    Flip disagreement & two backbone passes & no & \tbd & \tbd & \tbd & -- & -- \\
    Normalized disparity & channel 64 & no & \tbd & \tbd & \tbd & -- & -- \\
    Temporal spread & channel 65 & yes & \tbd & \tbd & \tbd & -- & -- \\
    \midrule
    Head, depth only (\texttt{m1-depth-only}) & channel 64 & no & \tbd & \tbd & \tbd & \tbd & \tbd \\
    Head, temporal only (\texttt{m1-temporal-only}) & channel 65 & yes & \tbd & \tbd & \tbd & \tbd & \tbd \\
    Head, features only (\texttt{m1-features-only}) & decoder features & no & \tbd & \tbd & \tbd & \tbd & \tbd \\
    Head, no temporal (\texttt{m1-no-temporal}) & all but channel 65 & no & \tbd & \tbd & \tbd & \tbd & \tbd \\
    Head, all inputs (\texttt{run-v1-radar}) & all 68 channels & yes & \tbd & \tbd & \tbd & \tbd & \tbd \\
    \midrule
    Oracle & true error & -- & 0 & \tbd & 1 & -- & -- \\
    \bottomrule
  \end{tabular}
  \caption{All methods on the nuScenes validation split (\tbd{} images, \tbd{} labeled cells), with the 95\% interval of each score from a bootstrap over the 150 scenes. ``Lidar at inference'' marks the methods that use the lidar fit of the depth scale (Section~\ref{sec:sensors}); a head so marked is an upper bound for a deployed system. The ECE is after the logit correction (Section~\ref{sec:protocol}). \todo{Fill from the per-image scores of every method on the same images, and add the intervals.}}
  \label{tab:main}
\end{table*}

\begin{figure}[t]
  \centering
  \placeholderfig{0.9\columnwidth}{4.2cm}{Sparsification curves. Mean AbsRel of the remaining cells against the fraction removed, for the oracle, random order, flip disagreement, the features-only head, and the full head. \texttt{plot.py}, \texttt{sparsification\_figure}.}
  \caption{\todo{Caption. Say what the axes are, which images the curves average over, and what the gap to the oracle means.}}
  \label{fig:sparsification}
\end{figure}

% ----------------------------------------------------------------------------
\subsection{Adverse Conditions}
\label{sec:conditions}
% MUST SHOW: the scores of the main methods for day, rain, and night, with
% intervals, next to the scene counts (111 / 24 / 15) and the label counts of
% Table 2. This is the claim in the title, so it gets its own table.
% Say which gaps between conditions the intervals support and which they do
% not: night has 15 scenes and rain 24.
% Also say whether the head keeps its ranking quality where the depth itself
% is worst, and whether its gain over flip disagreement grows or shrinks at night.

\todo{One bold sentence with the finding, then the numbers of Table~\ref{tab:conditions}.}

\begin{table*}[t]
  \centering
  \small
  \setlength{\tabcolsep}{4pt}
  \begin{tabular}{@{}l c c c c c c c c c@{}}
    \toprule
    & \multicolumn{3}{c}{\textbf{Day} (111 scenes)} & \multicolumn{3}{c}{\textbf{Rain} (24 scenes)} & \multicolumn{3}{c}{\textbf{Night} (15 scenes)} \\
    \cmidrule(lr){2-4} \cmidrule(lr){5-7} \cmidrule(l){8-10}
    \textbf{Method} & AUSE $\downarrow$ & AURG $\uparrow$ & AUROC $\uparrow$ & AUSE $\downarrow$ & AURG $\uparrow$ & AUROC $\uparrow$ & AUSE $\downarrow$ & AURG $\uparrow$ & AUROC $\uparrow$ \\
    \midrule
    Flip disagreement & \tbd & \tbd & \tbd & \tbd & \tbd & \tbd & \tbd & \tbd & \tbd \\
    Temporal spread & \tbd & \tbd & \tbd & \tbd & \tbd & \tbd & \tbd & \tbd & \tbd \\
    Head, features only & \tbd & \tbd & \tbd & \tbd & \tbd & \tbd & \tbd & \tbd & \tbd \\
    Head, all inputs & \tbd & \tbd & \tbd & \tbd & \tbd & \tbd & \tbd & \tbd & \tbd \\
    \bottomrule
  \end{tabular}
  \caption{Scores by condition on the validation split, with 95\% intervals from the bootstrap over the scenes of each group. Table~\ref{tab:labelcounts} gives the images and labels behind each column. \todo{Fill, with intervals; mark the differences between methods that the paired intervals support.}}
  \label{tab:conditions}
\end{table*}

% ----------------------------------------------------------------------------
\subsection{Which Inputs Matter}
\label{sec:input-ablation}
% MUST SHOW: what each input group adds, in AUSE and AURG on validation, with
% the interval of the difference from the baseline. The dev losses below are
% final and come from confhead/results/model_checkpoints.md (complete cache);
% the three findings are carried over from the midterm report with those
% losses. Restate each finding in AUSE and AURG once the runs are scored, and
% keep the dev loss only as a column.

Table~\ref{tab:input-ablation} lists the input ablation. The two runs with the same settings, \texttt{run-v1-radar} and its repeat, differ by 0.0008 in dev loss, which is the only measure of repeat noise on the dev loss. \todo{The findings below are stated on the dev loss, which shows that a run learned and does not rank the runs. Restate each with the AUSE, the AURG, and the paired interval of the difference from the baseline.}

\textbf{The decoder features carry most of the signal.} Without them the dev loss rises by 0.0540, the largest cost of removing any one input group, and the features alone come within 0.0094 of all inputs. A single channel on its own, depth or the temporal spread, is far worse than any run that reads the features.

\textbf{The temporal spread helps a little, and the depth channel less.} Removing the temporal spread costs 0.0059 and removing the depth channel 0.0018.

\textbf{Radar makes no difference on the dev loss.} The run without radar is within the repeat gap of the baseline, which is what the coverage of radar predicts: 95.2\% of the labeled cells have no radar cell on them \todo{Refresh on the complete cache; measured on the partial cache}. A tie over all labeled cells leaves open whether the head uses radar where radar exists. \todo{Compare \texttt{run-v1-radar} with \texttt{run-v0-no-radar} on the radar-valid cells alone (checklist B2; no retraining) and report the result here, as a number with its interval.}

\begin{table*}[t]
  \centering
  \small
  \setlength{\tabcolsep}{5pt}
  \begin{tabular}{@{}l l c c c c r r r r r@{}}
    \toprule
    \textbf{Run} & \textbf{What changes} & \textbf{Feat.} & \textbf{Depth} & \textbf{Temp.} & \textbf{Radar} & \textbf{Ch.} & \textbf{Dev loss} & \textbf{AUSE} $\downarrow$ & \textbf{AURG} $\uparrow$ & \textbf{$\Delta$AUSE} \\
    \midrule
    \texttt{run-v1-radar} & Nothing: the baseline & $\bullet$ & $\bullet$ & $\bullet$ & $\bullet$ & 68 & 0.5589 & \tbd & \tbd & 0 \\
    \texttt{run-v1-radar-repeat} & Nothing: a repeat & $\bullet$ & $\bullet$ & $\bullet$ & $\bullet$ & 68 & 0.5581 & \tbd & \tbd & \tbd \\
    \texttt{run-v0-no-radar} & No radar & $\bullet$ & $\bullet$ & $\bullet$ & -- & 66 & 0.5580 & \tbd & \tbd & \tbd \\
    \texttt{m1-no-depth} & No depth channel & $\bullet$ & -- & $\bullet$ & $\bullet$ & 67 & 0.5607 & \tbd & \tbd & \tbd \\
    \texttt{m1-no-temporal} & No temporal spread & $\bullet$ & $\bullet$ & -- & $\bullet$ & 67 & 0.5648 & \tbd & \tbd & \tbd \\
    \texttt{m1-features-only} & Decoder features alone & $\bullet$ & -- & -- & -- & 64 & 0.5683 & \tbd & \tbd & \tbd \\
    \texttt{m1-no-features} & No decoder features & -- & $\bullet$ & $\bullet$ & $\bullet$ & 4 & 0.6129 & \tbd & \tbd & \tbd \\
    \texttt{m1-depth-only} & Depth channel alone & -- & $\bullet$ & -- & -- & 1 & 0.6368 & \tbd & \tbd & \tbd \\
    \texttt{m1-temporal-only} & Temporal spread alone & -- & -- & $\bullet$ & -- & 1 & 0.6713 & \tbd & \tbd & \tbd \\
    \bottomrule
  \end{tabular}
  \caption{The input ablation: the input groups each run reads, its channel count, its lowest dev loss on the complete cache, and its scores on the validation split. $\Delta$AUSE is the difference from the baseline with its 95\% paired interval over scenes. \todo{Fill the three score columns.}}
  \label{tab:input-ablation}
\end{table*}

% ----------------------------------------------------------------------------
\subsection{Design Ablations}
\label{sec:design-ablation}
% MUST SHOW: that the conclusions do not hang on one setting. Three or four
% sentences and one compact table; the full list of runs is in Appendix B.
%   - radar settings (fit, radius, sweeps): does any of them make radar useful?
%   - tau at 0.05 and 0.2: these runs use other labels, so their dev loss is not
%     comparable. Compare them by AUSE and AURG on the AbsRel itself, which does
%     not depend on tau, or by AUROC at the evaluation tau of 0.10.
%   - label range up to 50 m, a quarter of the data, another seed.
% m2-radar-lidar-fit puts the lidar fit into the radar residual, so it is an
% upper bound like the full head (Section 3.5).

\todo{The finding in one bold sentence per group of runs, with the numbers of Table~\ref{tab:design-ablation}.}

\begin{table}[t]
  \centering
  \small
  \setlength{\tabcolsep}{3pt}
  \begin{tabular}{@{}l l r r r@{}}
    \toprule
    \textbf{Run} & \textbf{What changes} & \textbf{Dev loss} & \textbf{AUSE} $\downarrow$ & \textbf{AURG} $\uparrow$ \\
    \midrule
    \texttt{run-v1-radar} & Baseline & 0.5589 & \tbd & \tbd \\
    \midrule
    \texttt{m2-radar-lidar-fit} & Residual from lidar fit & 0.5534 & \tbd & \tbd \\
    \texttt{m2-radius0} & Radar radius 0 & 0.5591 & \tbd & \tbd \\
    \texttt{m2-radius2} & Radar radius 2 & 0.5581 & \tbd & \tbd \\
    \texttt{m2-sweeps1} & Keyframe sweep only & 0.5582 & \tbd & \tbd \\
    \midrule
    \texttt{m3-tau005} & $\tau = 0.05$ & 0.4362$^{\ast}$ & \tbd & \tbd \\
    \texttt{m3-tau02} & $\tau = 0.2$ & 0.6349$^{\ast}$ & \tbd & \tbd \\
    \texttt{m6-maxdepth50} & Labels up to 50\,m & 0.5633$^{\ast}$ & \tbd & \tbd \\
    \midrule
    \texttt{m6-quarter-data} & A quarter of the data & 0.5664 & \tbd & \tbd \\
    \texttt{m1-seed1} & Seed 1 & 0.5622$^{\dagger}$ & \tbd & \tbd \\
    \bottomrule
  \end{tabular}
  \caption{Design ablations, each one change from the baseline. $^{\ast}$Other labels, so the dev loss is not comparable with the baseline; the scores are computed on the AbsRel itself and are. $^{\dagger}$Other dev scenes. \todo{Fill the score columns.}}
  \label{tab:design-ablation}
\end{table}

% ----------------------------------------------------------------------------
\subsection{Calibration}
\label{sec:calibration}
% MUST SHOW: whether 1 - p can be read as a probability. A reliability diagram
% of the best head before and after the logit correction (subtract ln w+), with
% the ECE of each, overall and at night. If the head is still off after the
% correction, say in which direction (over- or under-confident).
% Optional: temperature scaling fitted on dev.

\todo{The finding, the two ECE values, and the reading of Figure~\ref{fig:reliability}.}

\begin{figure}[t]
  \centering
  \placeholderfig{0.9\columnwidth}{4.2cm}{Reliability diagram of the best head, before and after the logit correction. \texttt{plot.py}, \texttt{reliability\_figure}.}
  \caption{\todo{Caption.}}
  \label{fig:reliability}
\end{figure}

% ----------------------------------------------------------------------------
\subsection{Qualitative Results}
\label{sec:qualitative}
% MUST SHOW: what the confidence map looks like, where it is right, and where
% it fails. Rows: image, lidar AbsRel, confidence of the head, and flip
% disagreement for contrast. Columns: one day, one rain, one night frame where
% the head works, and at least one failure. Name the kind of surface in each
% (textureless road, object edge, sky, reflections, lens drops), since the
% project scope asks for a breakdown by surface and not only by weather.
% The map is dense and the labels are not: say what the head predicts where
% there is no lidar, and that those cells are not scored.

\todo{Two or three sentences on what Figure~\ref{fig:qualitative} shows, with the failure case named.}

\begin{figure*}[t]
  \centering
  \placeholderfig{0.95\textwidth}{6cm}{Qualitative grid. Image, lidar AbsRel, head confidence, and flip disagreement for day, rain, and night frames, with one failure case. \texttt{plot.py}, \texttt{qualitative\_figure}.}
  \caption{\todo{Caption. Name the scene and camera of each column and the color scale.}}
  \label{fig:qualitative}
\end{figure*}

% ----------------------------------------------------------------------------
\subsection{Downstream Use}
\label{sec:downstream}
% OPTIONAL. Cut this subsection if no number is ready by the results freeze,
% and move its idea to Section 7 as future work.
% MUST SHOW, if kept: one number that improves when the confidence is used.
% Candidates (checklist B4): (a) a frame selector that ranks whole images by
% mean p_error, with the error of the remaining images as the worst are
% dropped; (b) the radar scale-and-shift fit with each point weighted by the
% confidence of its cell; (c) the outlier rate at a fixed coverage after the
% least confident cells are dropped. Use the confidence as a weight or a gate;
% multiplying a depth by a probability only shrinks it.

\todo{The experiment, its number against the same pipeline without the confidence, and one sentence on what it shows. Cut if not ready.}
